% Part: normal-modal-logic
% Chapter: filtrations
% Section: S5-decidable

\documentclass[../../../include/open-logic-section]{subfiles}

\begin{document}

\olfileid{nml}{fil}{dec}

\olsection{\Log{S5} تصمیم‌پذیر است}
خاصیت مدل متناهی راهی آسان برای نشان‌دادن این امر به دست می‌دهد که
دستگاه‌های منطق موجهاتیِ داده‌شده به‌وسیلهٔ طرح‌واره‌ها \emph{تصمیم‌پذیر}
هستند (یعنی روشی محاسبه‌پذیر وجود دارد تا تعیین کند آیا !!a{formula}
در دستگاه !!{derivable} است یا نه).

\begin{thm}
  \Log{S5} تصمیم‌پذیر است.
\end{thm}

\begin{proof}
  بگذارید~$!A$ داده شده باشد و فرض کنید !!{propositional variable}های
  ظاهرشونده در~$!A$ در میان~$p_1$، \dots، $p_k$ باشند. چون برای هر
  $n$ تنها شمار متناهی‌ای مدل همگانیِ متناهی با~$n$ جهان وجود دارد که
  به~$p_1$، \dots، $p_k$ ارزش می‌دهند، می‌توانیم \emph{به‌طور موازی}
  همهٔ قضایای~\Log{S5} را با تولید نظام‌مند اثبات‌ها فهرست کنیم؛ و
  همهٔ مدل‌های همگانی متناهیِ دارای~$1$، $2$، \dots{} جهان را نیز
  فهرست کنیم و بیازماییم که آیا~$!A$ در جهانی از یکی از آن‌ها شکست
  می‌خورد. سرانجام یکی از این دو فرایند موازی پاسخ خواهد داد، زیرا بنا
  بر \olref[com][fra]{thm:generaldet} و \olref[fmp]{cor:S5fmp}، یا
  $!A$ !!{derivable} است یا در یک مدل همگانی متناهی شکست می‌خورد.
\end{proof}

اثبات بالا برای~\Log{S5} کار می‌کند، زیرا تصفیه‌های مدل‌های همگانی
خودبه‌خود همگانی‌اند. همین امر دربارهٔ بازتابی‌بودن و سریال‌بودن برقرار
است، اما برای خواص دیگر به کار بیشتری نیاز است.

\end{document}
